\chapter{Inversion}\label{geometry:ch:inversion}

\begin{center}\itshape
Four circles to the kissing come.\\
The smaller are the benter.\\
The bend is just the inverse of\\
The distance from the center.
\upshape(Frederick Soddy, \emph{The Kiss Precise}, Nature, 1936)
\end{center}

A circle can become a straight line if we change how we look at the plane. Inversion does this while keeping the sizes of angles. A ring of tangent circles can become a ring of equal circles, and a curved chain can become a stack between parallel lines. The simpler picture lets us calculate before bringing the answer back.

Descartes sent Elisabeth of Bohemia a relation for four tangent circles in November 1643. Soddy rediscovered it nearly three centuries later and published it as a poem in \emph{Nature} in 1936. The bend in the epigraph is curvature, the reciprocal radius. An enclosing circle receives a negative bend; a limiting straight line receives bend zero. With those signs, the relation is
\[
(k_1+k_2+k_3+k_4)^2=2(k_1^2+k_2^2+k_3^2+k_4^2).
\]
The centre coordinates satisfy related identities that will tell us what happens when one circle is replaced by another.

The prerequisites are complex numbers, determinants and the Euclidean transformations of Appendix~A. The M\"obius maps studied here have already appeared in Chapter~9's hyperbolic geometry. Inversion also has a physical use: Thomson's 1845 electrostatic construction replaces a spherical boundary by an image charge. One of the new exercises verifies that construction directly.

\begin{think}
Place a circle strictly inside another. Fit a circle in the gap touching both, then keep fitting the next tangent circle in one direction around the gap. If one starting choice closes after seven steps, will every starting choice do so? Reversing immediately onto the previous circle is not a new step around the gap.
\end{think}

\section{Inversion in a circle}

\begin{definition}\label{geometry:def:inversion}
For a circle with centre $O$ and radius $r>0$, its \emph{inversion} sends $P\ne O$ to $P'$ on the same ray with $OP\,OP'=r^2$. With $O=0$ in complex coordinates, it is $z\mapsto r^2/\bar z$.
\end{definition}

Applying the map twice returns $P$. It fixes its defining circle pointwise and exchanges inside with outside. Add a point $\infty$ and exchange it with $O$; the resulting extended plane is $\hat\C=\C\cup\{\infty\}$. The local coordinate $1/z$ near infinity is compatible with the sphere description from stereographic projection. Lines will include that point, so they can be treated alongside circles.

\begin{proposition}[distances]\label{geometry:prop:invdist}
For $P,Q\ne O$, the inverted distance is
\[
P'Q'=\frac{r^2\,PQ}{OP\,OQ}.
\]
\end{proposition}
\begin{proof}
Writing the points as nonzero $z,w$ gives
$|r^2/\bar z-r^2/\bar w|=r^2|z-w|/(|z||w|)$, including the case $P=Q$.
\end{proof}

A \emph{generalized circle} means an ordinary positive-radius circle or a line, with the line completed by infinity. In finite coordinates its equation is
\begin{equation}\label{geometry:eq:gencircle}
\begin{gathered}
A|z|^2+Bz+\bar B\bar z+C=0,\\
A,C\in\R,\quad B\in\C,\quad |B|^2>AC.
\end{gathered}
\end{equation}
For $A\ne0$, completing the square gives centre $-\bar B/A$ and radius $\sqrt{|B|^2-AC}/|A|$. For $A=0$, the strict inequality makes $B\ne0$ and the equation is a line. The set passes through zero exactly when $C=0$.

\begin{theorem}[inversion and circles]\label{geometry:thm:invcircles}
Inversion carries generalized circles to generalized circles. Lines through its centre map to themselves; other lines and circles through the centre are exchanged; circles avoiding the centre map to circles avoiding it. At an intersection of regular curves, inversion preserves the angle's magnitude and reverses its orientation.
\end{theorem}
\begin{proof}
Translate and scale to $O=0,r=1$. Put $w=1/\bar z$, so $z=1/\bar w$. Substitution into \eqref{geometry:eq:gencircle}, followed by multiplication by $|w|^2$, gives
\[
C|w|^2+Bw+\bar B\bar w+A=0.
\]
The coefficient roles of $A,C$ are exchanged and $|B|^2>AC$ is unchanged. Reading whether each coefficient vanishes gives the listed circle/line cases. The points zero and infinity complete the images at the omitted finite coordinates.

At finite $z\ne0$, a tangent vector $v$ transforms to $-\bar v/\bar z^2$. Complex conjugation reverses directed angles, while multiplication by the nonzero factor preserves their magnitude and orientation. Their combination gives the claim. At zero or infinity, use the local reciprocal coordinate on the infinite side: the expression is again conjugation followed by a nonzero similarity. The angle statement concerns curves with nonzero well-defined tangents in these charts.
\end{proof}

Two distinct generalized circles are tangent when they have one intersection on the extended plane; parallel lines meet only at infinity. A bijective inversion preserves that condition. Ordinary circles of radii $\rho_1,\rho_2$ whose centres are distance $d$ apart are orthogonal precisely when $d^2=\rho_1^2+\rho_2^2$, by the right triangle formed at an intersection.

\begin{proposition}\label{geometry:prop:orthfixed}
Inversion in $\omega$ preserves each circle orthogonal to $\omega$. Conversely, any circle through a pair of distinct inverse points $P,P'$ is orthogonal to $\omega$.
\end{proposition}
\begin{proof}
For a circle $\gamma$ with centre $C$ and radius $\rho$, the \emph{power} of $O$ is $OC^2-\rho^2$. Along a line $O+tu$ with $|u|=1$, intersection parameters solve
\[
t^2-2\langle u,C-O\rangle t+OC^2-\rho^2=0.
\]
Their product is the power. If it is positive and the line meets the circle, both intersections are on the same ray.

If $\gamma$ is orthogonal to the inversion circle of radius $r$, the right-triangle criterion gives power $r^2$. Each radial intersection pair is therefore an inverse pair; at a tangency the coincident pair is fixed. Thus the circle maps onto itself. Conversely a circle through distinct $P,P'$ has power $OP\,OP'=r^2$, because the two points are on the same ray. Its centre distance satisfies $OC^2=\rho^2+r^2$, giving orthogonality.
\end{proof}

\section{Ptolemy's theorem}

Ptolemy used circle chords to build trigonometric tables around the second century. Inversion turns the relation between four points into the triangle inequality for three.

\begin{theorem}[Ptolemy's inequality]\label{geometry:thm:ptolemy}
For any four plane points,
\[
AB\,CD+AD\,BC\ge AC\,BD.
\]
Equality holds if two adjacent named points coincide, meaning $A=B$, $B=C$, $C=D$ or $D=A$; otherwise it holds exactly when all four are distinct and lie in cyclic order $A,B,C,D$ on a generalized circle. Either orientation is allowed; a line is viewed as its circle completed by infinity.
\end{theorem}
\begin{proof}
An adjacent coincidence gives equality by substitution. If there is an opposite coincidence and no adjacent one, the right side is zero and the left side is strictly positive. It remains to consider four distinct points.

Invert with radius one about $A$. The triangle inequality $B'C'+C'D'\ge B'D'$ becomes
\[
\frac{BC}{AB\,AC}+\frac{CD}{AC\,AD}\ge\frac{BD}{AB\,AD}.
\]
Multiplying by the positive product $AB\,AC\,AD$ gives the claim. Equality means $C'$ lies on the segment from $B'$ to $D'$. The inverse line is a generalized circle through $A$, with $B,C,D$ in the stated order away from $A$. Conversely that cyclic order puts $C'$ between the other two, giving equality.
\end{proof}

For a cyclic quadrilateral the products of opposite sides add to the product of diagonals. In a rectangle this is Pythagoras. Choosing chords in a unit-diameter circle also gives the sine addition formula: inversion has recovered a trigonometric relation from a straight-line inequality.

\section{Steiner's porism}

For two ordinary circles, define their \emph{inversive distance} by
\[
\delta=\frac{\rho_1^2+\rho_2^2-d^2}{2\rho_1\rho_2}.
\]
Internal tangency gives $1$, external tangency gives $-1$, and orthogonality gives zero. When the circles cross, the cosine rule identifies this with the cosine of the angle between their radius vectors at an intersection.

\begin{proposition}[inversive distance is invariant]\label{geometry:prop:invdistance}
An inversion mapping two circles avoiding its centre to ordinary circles preserves $|\delta|$.
\end{proposition}
\begin{proof}
For centre $c$ and radius $\rho$, normalize the equation coefficients to
$X=(A,B,C)=\rho^{-1}(1,-\bar c,|c|^2-\rho^2)$. On their real coefficient space put
\[
\langle X,Y\rangle=\operatorname{Re}(B_X\bar B_Y)
-\tfrac12(A_XC_Y+A_YC_X).
\]
Then $\langle X,X\rangle=1$, and expansion gives
\[
\langle X_1,X_2\rangle
=\frac{\rho_1^2+\rho_2^2-|c_1-c_2|^2}{2\rho_1\rho_2}=\delta.
\]
Similarities preserve this distance ratio directly. It suffices therefore to check unit inversion about zero. It exchanges $A,C$, preserving the displayed bilinear form. The resulting coefficient vector describes the image circle, but its normalization to the radius convention may differ by a real factor. Both normalized vectors have squared norm one, so that factor is $1$ or $-1$. Inner products can change sign but not absolute value, proving the claim.
\end{proof}

The possible sign change records whether inversion exchanges a circle's inside and outside. Keeping $|\delta|$ avoids choosing an orientation for each of the ordinary circles.

\begin{lemma}[concentric position]\label{geometry:lem:concentric}
Any two disjoint ordinary circles can be carried by an inversion to concentric circles.
\end{lemma}
\begin{proof}
Already concentric circles stay concentric under inversion about their common centre. Otherwise put their centres at $0,d$ with $d>0$ and radii $\rho_1,\rho_2$. Seek real $p,q$ inverse with respect to both:
\[
pq=\rho_1^2,\qquad(p-d)(q-d)=\rho_2^2.
\]
They are roots of $x^2-sx+\rho_1^2$, with $s=(\rho_1^2+d^2-\rho_2^2)/d$. Its discriminant is
\[
s^2-4\rho_1^2=
\frac{((d-\rho_1)^2-\rho_2^2)((d+\rho_1)^2-\rho_2^2)}{d^2}>0.
\]
The sign follows from external separation $d>\rho_1+\rho_2$ or strict nesting $d<|\rho_1-\rho_2|$. Thus the roots are distinct, nonzero and on the same side of zero, one with absolute value less than $\rho_1$. Neither lies on either circle: an inverse point on its inversion circle is its own inverse, contrary to $p\ne q$.

Invert about $p$. Choose two distinct circles through $p,q$. By the preceding orthogonality proposition, each is orthogonal to both given circles. Their images are distinct lines through $q'$. The two given image circles remain ordinary because they avoid $p$, and are orthogonal to both image lines. A line orthogonal to a circle passes through its centre, so both centres are the intersection $q'$. The images are concentric.
\end{proof}

\begin{theorem}[Steiner's porism]\label{geometry:thm:steiner}
Let $\Gamma_1$ lie strictly inside $\Gamma_2$. Consider circles whose discs lie inside $\Gamma_2$ and outside the interior of $\Gamma_1$, tangent to both boundaries, advancing consistently in one direction from each circle to its next tangent neighbour, without reversing onto the previous circle. If the first return closes after $n$ distinct circles for one starting choice, it closes after the same $n$ circles and number of turns for every starting choice.
\end{theorem}
\begin{proof}
Use the concentric-position construction with the inner circle as the first circle. For nonconcentric nested circles its root closest to the inner centre lies strictly inside that circle. Choosing this inversion centre keeps it outside the gap and every chain circle. The gap maps homeomorphically to a bounded concentric annulus; for already concentric circles invert about their common centre. Write its boundary radii as $a<b$.

Every circle in that annulus tangent to both boundaries has radius $(b-a)/2$ and centre at distance $(a+b)/2$ from their common centre. Consecutive centres have chord length $b-a$ and therefore angular separation
\[
\theta=2\arcsin\frac{b-a}{b+a},\qquad 0<\theta<\pi.
\]
The chosen forward branch keeps the same sign of this step. A first return after $n$ steps is exactly $n\theta=2\pi k$, with coprime positive integers $n,k$. It does not involve the starting angle. Rotating the whole concentric chain reaches every possible first circle, and inverting back gives the result. The annular homeomorphism preserves the absolute winding count $k$; its sign depends on orientation. In particular $n\ge3$.
\end{proof}

\begin{figure}[ht]
\centering
\fig{geometry-ch03-steiner}
\caption{A Steiner chain of seven circles. It is the image, under an inversion, of a ring of seven equal circles between two concentric
circles, which is why it closes up wherever it starts.}
\label{geometry:fig:steiner}
\end{figure}

A \emph{porism} says that the existence of one member forces a family of solutions. Here the family comes from freely rotating the first circle in the concentric picture. Appendix~D will meet another porism for polygons between conics.

\section{The Pappus chain}

Choose $C$ strictly between $A,B$. The \emph{arbelos}, or shoemaker's knife, lies inside the upper half-disc on $AB$ and outside the upper half-discs on $AC,CB$. The two smaller radii are positive. Pappus studied a chain of inscribed circles advancing up this gap from the boundary circle on $CB$.

\begin{theorem}[Pappus]\label{geometry:thm:pappus}
Let $C_0$ be the circle with diameter $CB$. Let $C_1,C_2,\ldots$ be the successive circles in the upper arbelos tangent to the circles on $AB,AC$ and to the preceding member, advancing without backtracking. If $d_n$ is the diameter of $C_n$, its centre is at height $n d_n$ above $AB$.
\end{theorem}
\begin{proof}
For a fixed $n$, the point $A$ is outside $C_n$, so its power relative to $C_n$ is positive. Invert about $A$ with radius equal to the square root of that power. The orthogonality proposition makes $C_n$ invariant as a circle.

The circles on $AB,AC$ pass through $A$ and have centres on $AB$. They become parallel lines perpendicular to $AB$: the baseline remains itself and its right-angle intersections remain right angles. The chain circles avoid $A$ and become ordinary circles tangent to both parallel lines, so all have the same diameter. Inversion about a point on $AB$ preserves its upper half-plane and commutes with reflection in $AB$. Thus $C_0$'s image has centre on $AB$, and the advancing chain becomes a vertical stack whose $j$th centre is $j$ diameters above that line. The image of $C_n$ is the original circle itself, with the same centre and diameter, giving its height $nd_n$. This uses invariance of the circle, not a claim that inversion fixes its centre point. For $n=0$ the statement also follows immediately from the original baseline centre.
\end{proof}

\begin{figure}[ht]
\centering
\fig{geometry-ch03-pappus}
\caption{The Pappus chain in an arbelos. The centre of $C_n$ is exactly $n$ diameters of $C_n$ above the line $AB$ (rust), as the computation
from the inverted picture confirms.}
\label{geometry:fig:pappus}
\end{figure}

\section{Descartes' circle theorem}

Orient a proper tangent-circle configuration by giving ordinary circles positive curvature $1/\rho$ and an enclosing circle negative curvature $-1/\rho$. The enclosing member, if present, contains the other discs; otherwise the discs are mutually external. Distinct tangencies exclude the configurations where all circles merely touch at the same point. Use the signed coefficient vector
\[
X=k(1,-\bar c,|c|^2-\rho^2).
\]
It has $\langle X,X\rangle=1$, and tangent pairs in this orientation have inner product $-1$, by the earlier inversive-distance calculation. The first coordinate is the signed curvature $k$, while $-\bar B=kc$ records the centre weighted by curvature.

\begin{think}
Three externally tangent coins have radii $1,2,3$. Which radius should the circle enclosing all three have? The quadratic below gives a negative curvature for that outer solution and a positive curvature for the small inner gap.
\end{think}

\begin{theorem}[Descartes 1643, Soddy 1936; Lagarias, Mallows and Wilks 2002]\label{geometry:thm:descartes}
For four circles pairwise tangent at six distinct points, with the signed curvatures and centres just described,
\[
\left(\sum_i k_i\right)^2=2\sum_i k_i^2,\qquad
\left(\sum_i k_ic_i\right)^2=2\sum_i(k_ic_i)^2,
\qquad
\left(\sum_i k_i\right)\left(\sum_i k_ic_i\right)=2\sum_i k_i(k_ic_i).
\]
\end{theorem}
\begin{proof}
In real coordinates $(A,\operatorname{Re}B,\operatorname{Im}B,C)$, the coefficient form and its inverse are
\[
\Omega=\begin{pmatrix}0&0&0&-1/2\\0&1&0&0\\0&0&1&0\\-1/2&0&0&0\end{pmatrix},
\qquad
\Omega^{-1}=\begin{pmatrix}0&0&0&-2\\0&1&0&0\\0&0&1&0\\-2&0&0&0\end{pmatrix}.
\]
Let $M$ have the four vectors as rows. Their pairings give
$M\Omega M^\top=G=2I-J$, with $J$ the all-ones $4\times4$ matrix. Since $J^2=4J$, direct multiplication gives $G^{-1}=I/2-J/4$. Invertibility of $G$ makes $M$ invertible, and hence
\[
\Omega^{-1}=M^\top G^{-1}M.
\]
For real or complex column vectors $u,v$ this implies
$u^\top\Omega^{-1}v=(Mu)^\top G^{-1}(Mv)$. Choose first $u=v=e_1$. Its left side is zero and $Me_1=(k_i)$, giving
$0=\tfrac12\sum k_i^2-\tfrac14(\sum k_i)^2$.

Next choose $u=v=\phi=(0,-1,i,0)^\top$. The left side is $(-1)^2+i^2=0$, while $(M\phi)_j=-\operatorname{Re}B_j+i\operatorname{Im}B_j=-\bar B_j=k_jc_j$. This gives the second identity. Finally take $u=e_1,v=\phi$; the left side is again zero and gives the third identity. These are bilinear complex calculations, without conjugation in the transpose products.
\end{proof}

A normalized coefficient vector with $A=0$ represents an oriented line: then $|B|=1$ and its equation is $Bz+\bar B\bar z+C=0$. Its curvature is zero. Its coefficient $-\bar B$ still makes sense, but it is not curvature times a finite centre. This distinction matters when a replacement circle becomes a line.

\begin{proposition}[the swap rule]\label{geometry:prop:swap}
Let $X_1,X_2,X_3,X_4$ be an oriented Descartes configuration. The other oriented generalized circle tangent to the first three has vector $X_4'$ satisfying
\[
X_4+X_4'=2(X_1+X_2+X_3).
\]
The curvature identity holds also for a line replacement; when both fourth members are ordinary circles, the centre identity is
\begin{equation}\label{geometry:eq:descartesswap}
\begin{aligned}
k_4+k_4'&=2(k_1+k_2+k_3),\\
k_4c_4+k_4'c_4'&=2(k_1c_1+k_2c_2+k_3c_3).
\end{aligned}
\end{equation}
For a line, replace each centre-weighted term by its defined coefficient $w=-\bar B$ rather than dividing by zero curvature.
\end{proposition}
\begin{proof}
Put $Y=X_1+X_2+X_3$. Its pairings with the first three vectors are $-1$, and their Gram matrix $G_3=2I-J$ has determinant $-4$, so they are independent and span a nondegenerate three-dimensional subspace. All vectors with those same three pairings form an affine line $Y+\R N$, where $N\ne0$ is orthogonal to that span.

We need the sign of $\langle N,N\rangle$. Form the basis matrix $L$ with columns $X_1,X_2,X_3,N$. It is invertible because a vector in their span and orthogonal to it would be zero. Taking determinants of its Gram matrix gives
\[
-4\langle N,N\rangle
=\det(L)^2\det\Omega=-\tfrac14\det(L)^2,
\]
so $\langle N,N\rangle>0$. This proves the required sign directly. Also $\langle Y,Y\rangle=3-6=-3$ and $\langle Y,N\rangle=0$. Thus norm-one points on the affine line satisfy
$t^2\langle N,N\rangle=4$ and are the two distinct vectors $Y\pm t_0N$. A norm-one coefficient vector describes an oriented ordinary circle when $A\ne0$, or a line when $A=0$. These are the two possible replacements, so their sum is $2Y$. Reading their first and $-\bar B$ coordinates gives the scalar identities, with the stated line convention.
\end{proof}

For the coins $1,2,3$, the positive curvatures are $1,1/2,1/3$. The quadratic gives
\[
k_4=\frac{11}{6}\pm2\sqrt{\frac12+\frac16+\frac13}
=\frac{11}{6}\pm2.
\]
The negative root $-1/6$ means an enclosing circle of radius six, and the positive root $23/6$ gives the inner radius $6/23$. Exercise~\ref{ex:v9-descartes-line} treats a configuration whose other root is zero and identifies the actual line.

\section{Apollonian gaskets}

Start with four tangent circles, one enclosing the other three. Fill each curvilinear triangular gap by the circle tangent to its three boundary arcs, and repeat on the new gaps. This gives an \emph{Apollonian gasket}. The residual set has zero planar area and Hausdorff dimension approximately $1.3057$. Those are quoted results about the infinite construction, not consequences proved by drawing a finite stage; McMullen's work on dimension is listed below.

\begin{corollary}[integral gaskets]\label{geometry:cor:integral}
Integer starting curvatures give integer curvatures at every stage. If the initial weighted centres $k_ic_i$ are Gaussian integers as well, every later weighted centre is a Gaussian integer.
\end{corollary}
\begin{proof}
A new circle is the other member of a Descartes quadruple. The swap formulas replace its curvature by twice the sum of the other three minus the old curvature, and do the same to the weighted centre. Both operations preserve the respective integer rings. Induction over finite generations proves the assertion for every circle added. In this bounded packing all new members are ordinary finite circles inside gaps.
\end{proof}

\begin{figure}[ht]
\centering
\fig{geometry-ch03-gasket}
\caption{The Apollonian gasket with root quadruple $(-1,2,2,3)$: the outer circle has radius $1$, the two large circles radius $\frac12$, and every
circle has an integer curvature. The circles of curvature $6$ come from $2\cdot(-1+2+3)-2$.}
\label{geometry:fig:gasket}
\end{figure}
\FloatBarrier

Which integers occur is subtler than their integrality. For primitive integral gaskets, Bourgain and Kontorovich proved a density-one result in 2014 for integers passing the local congruence tests. Haag, Kertzer, Rickards and Stange's 2023 work, published in 2024, exhibited further quadratic-reciprocity obstructions to a full local-global conjecture. These research results are quoted; the exercises use only the elementary swap rule and its residue classes.

\section{M\"obius transformations}

A \emph{M\"obius transformation} is
\[
f(z)=\frac{az+b}{cz+d},\qquad ad-bc\ne0,
\]
on the extended complex plane. If $c\ne0$, its finite pole $-d/c$ maps to infinity and infinity maps to $a/c$. If $c=0$, it is an affine complex map and fixes infinity. Matrix multiplication describes composition, and inverse matrices describe inverses. Nonzero scalar multiples give the same map. Conversely a matrix inducing the identity has $cz^2+(d-a)z-b=0$ for every finite $z$, forcing $c=b=0,a=d$; thus scalar matrices are exactly the kernel of this matrix description.

\begin{theorem}[M\"obius maps]\label{geometry:thm:mobius}
\begin{enumerate}
\item Every M\"obius map is a composition of translations, nonzero complex multiplications and $z\mapsto1/z$. It preserves generalized circles and directed angles.
\item For any two ordered triples of distinct extended-plane points, there is exactly one M\"obius map carrying the first triple to the second.
\item For four distinct extended-plane points, the cross-ratio
\[
[z_1,z_2;z_3,z_4]=
\frac{(z_1-z_3)(z_2-z_4)}{(z_1-z_4)(z_2-z_3)}
\]
is invariant under M\"obius maps. Values with one point at infinity are defined by limits. The cross-ratio is real exactly when the points lie on a generalized circle.
\end{enumerate}
\end{theorem}
\begin{proof}
(1) For $c=0$ the map is a nonzero multiplication followed by translation. For $c\ne0$, write
\[
f(z)=\frac ac-\frac{ad-bc}{c(cz+d)}.
\]
This gives the stated decomposition. Reciprocal $1/z$ is circle inversion followed by conjugation. Both reverse orientation, so their composition preserves it; translations and nonzero multiplications preserve it too. Their circle properties, including infinity, prove the claim.

(2) For finite distinct $z_1,z_2,z_3$, the map
\[
g(z)=\frac{(z-z_1)(z_2-z_3)}{(z-z_3)(z_2-z_1)}
\]
sends them to $0,1,\infty$. If $z_1,z_2$ or $z_3$ is infinity, the respective limit formulas are
$(z_2-z_3)/(z-z_3)$, $(z-z_1)/(z-z_3)$, and $(z-z_1)/(z_2-z_1)$.
Construct $h$ similarly for the target triple; then $h^{-1}g$ works. For uniqueness, normalize both triples this way. A map fixing $\infty,0,1$ has $c=0,b=0,a=d$, so is the identity.

(3) Translations and multiplications cancel from the ratio. For reciprocation, the difference identity $1/z_i-1/z_j=(z_j-z_i)/(z_iz_j)$ makes all extra factors cancel. Values involving zero or infinity follow by continuity of the limit formulas. Finally normalize $z_2,z_3,z_4$ to $1,0,\infty$. The ratio becomes $[f(z_1),1;0,\infty]=f(z_1)$. The unique generalized circle through $1,0,\infty$ is the real line with infinity, so the circle-preservation property gives the real-value criterion in both directions.
\end{proof}

\begin{remark}
Separate the identity before classifying other M\"obius maps. Normalize a representing matrix to determinant one; its two possible signs give the same squared trace. For a nonidentity map, squared trace in $[0,4)$ is the elliptic case, conjugate to a nontrivial unit rotation; squared trace four is parabolic, conjugate to $z+1$; all remaining squared traces are loxodromic, conjugate to $\lambda z$ with $|\lambda|\ne1$. Exercise~\ref{geometry:ex:trace} proves these descriptions through fixed points. The real determinant-one maps preserving the upper half-plane give its orientation-preserving hyperbolic isometries, as in Chapter~9; the full isometry group also has orientation-reversing maps.
\end{remark}

\section{A gem: Feuerbach's theorem}

The incircle touches the three sides from inside the triangle. Each excircle lies outside, touching one side and the extensions of the other two; their centres use one internal and two external angle bisectors. The nine-point circle passes through the three side midpoints, the three altitude feet and the three midpoints between vertices and orthocentre. Feuerbach's theorem relates it to the incircle and excircles. Coxeter and Greitzer give an inversion proof; we will use unit-circle coordinates so that the exceptional equilateral case stays visible.

\begin{proposition}[three tests]\label{geometry:prop:tests}
For complex $a\ne b$:
\begin{enumerate}
\item $z$ lies on line $ab$ exactly when $(z-a)/(b-a)$ is real.
\item For $z\ne a$, line $az$ is perpendicular to $ab$ exactly when this ratio is purely imaginary.
\item For any $c$, the counterclockwise signed area is
\[
[abc]=\tfrac12\operatorname{Im}(\overline{(b-a)}(c-a))
=\frac i4\begin{vmatrix}a&\bar a&1\\b&\bar b&1\\c&\bar c&1\end{vmatrix}.
\]
\end{enumerate}
\end{proposition}
\begin{proof}
For $z\ne a$, the quotient's argument is the directed angle from $b-a$ to $z-a$. Real values give parallel vectors and purely imaginary values give perpendicular ones. At $z=a$, the first statement holds directly because the quotient is zero; the second excludes this undefined line. For the area write $u=b-a$, $v=c-a$. Its real determinant is $u_1v_2-u_2v_1=\operatorname{Im}(\bar uv)$. Subtracting the first row from the others in the complex determinant leaves $u\bar v-\bar uv=-2i\operatorname{Im}(\bar uv)$, so multiplication by $i/4$ gives the same half-determinant, also for collinear points.
\end{proof}

\begin{lemma}[chords of the unit circle]\label{geometry:lem:chord}
If $a\ne b$ and $|a|=|b|=1$, then:
\begin{enumerate}
\item line $ab$ has equation $z+ab\bar z=a+b$;
\item the perpendicular foot from $z$ to this line is
$f=(a+b+z-ab\bar z)/2$, and reflection in the line is $z\mapsto a+b-ab\bar z$.
\end{enumerate}
\end{lemma}
\begin{proof}
A real quotient equals its conjugate. Using $\bar a=1/a$, $\bar b=1/b$, that equality becomes
$z-a=b-ab\bar z$, giving (1). For the stated $f$, conjugation gives
$ab\bar f=(a+b+ab\bar z-z)/2$, so $f+ab\bar f=a+b$ and $f$ lies on the line. Put $w=a+b-z-ab\bar z=2(f-z)$. Then $\bar w=w/(ab)$ and $\overline{b-a}=(a-b)/(ab)$, making the quotient $w/(b-a)$ equal to minus its conjugate. If $f\ne z$ this is the perpendicular direction; if $f=z$, the point was already on the line. The reflected point is $2f-z$, giving the formula.
\end{proof}

\begin{proposition}[the incentre]\label{geometry:prop:incentre}
For a nondegenerate triangle with distinct vertices $a,b,c$ on the unit circle, square roots $u,v,w$ can be chosen so that the arc midpoints opposite the respective vertices are $-vw,-wu,-uv$. With this choice the incentre coordinate is $\iota=-(uv+vw+wu)$.
\end{proposition}
\begin{proof}
Read the vertices counterclockwise, interchanging labels if necessary, and choose arguments $a=e^{2i\alpha}$, $b=e^{2i\beta}$, $c=e^{2i\gamma}$ with $\alpha<\beta<\gamma<\alpha+\pi$. The arc midpoints of $AB,BC$ avoiding the third vertex are $e^{i(\alpha+\beta)},e^{i(\beta+\gamma)}$. The corresponding $CA$ midpoint is $-e^{i(\gamma+\alpha)}$. Taking $u=e^{i\alpha}$, $v=-e^{i\beta}$, $w=e^{i\gamma}$ gives the stated three signs.

The internal angle bisector from $a=u^2$ passes through $-vw$, by equality of the two arcs it cuts off. Their chord line has equation $z-u^2vw\bar z=u^2-vw$. For $\iota=-(uv+vw+wu)$ we have $\bar\iota=-(u+v+w)/(uvw)$, so
\[
\iota-u^2vw\bar\iota
=-(uv+vw+wu)+u(u+v+w)=u^2-vw.
\]
Thus $\iota$ lies on that internal bisector line, and the symmetric computations put it on the other two. Their intersection is the incentre, proving the formula.
\end{proof}

Appendix~A gives orthocentre $h=a+b+c$ and circumcentre zero. The circle with centre $n=h/2$ and radius $1/2$ passes through all nine named points. For example, side midpoint $(a+b)/2$ differs from $n$ by $-c/2$; the altitude foot from $a$, computed by the chord lemma, is $(h-bc/a)/2$, differing by $-bc/(2a)$; the midpoint of $a,h$ differs by $a/2$. Each difference has norm $1/2$, and the other six points follow cyclically. This verifies the nine-point circle used in the theorem.

\begin{theorem}[Feuerbach, 1822]\label{geometry:thm:feuerbach}
For a nondegenerate nonequilateral triangle, the nine-point circle is internally tangent to the incircle and externally tangent to all three excircles. In the equilateral case, the nine-point circle and incircle coincide, while the three external excircle tangencies still hold.
\end{theorem}
\begin{proof}
Scale to circumradius $R=1$ and use the chosen roots. Write $s_1=u+v+w$, $s_2=uv+vw+wu$, $s_3=uvw$. Since the roots have modulus one, $\bar s_1=s_2/s_3$. The incentre is $\iota=-s_2$ and the nine-point centre is $n=(s_1^2-2s_2)/2$. Hence
\[
n-\iota=\tfrac12s_1^2,\qquad NI=\tfrac12|s_1|^2.
\]
The chord-foot formula gives the inradius as
$r=|b+c-\iota-bc\bar\iota|/2$. Expansion yields
\[
u(b+c-\iota-bc\bar\iota)
=(u+v)(v+w)(w+u)=s_1s_2-s_3
=s_3(|s_1|^2-1).
\]
Thus $r=|1-|s_1|^2|/2$. Also $|\iota|=|s_2|=|s_1|$. The incentre lies inside the triangle, hence strictly inside its circumdisc, so $|s_1|<1$ and
\[
NI=\tfrac12|s_1|^2=\tfrac12-r.
\]
If $NI>0$ this is precisely internal tangency of radii $1/2,r$. If $NI=0$, the two centres and radii agree, giving coincident circles. This happens exactly when the incentre is the circumcentre. Equal distances from that centre to the three side chords give equal chord lengths, so the triangle is equilateral; the converse is immediate. Rescaling also gives Euler's identity $OI^2=R(R-2r)$.

For the $A$-excentre change $u$ to $-u$. The unchanged squares still represent the same vertices, and the resulting coordinate $\iota_A=-(-uv+vw-wu)$ lies on the internal bisector at $A$ and the external bisectors at $B,C$: the latter now use the opposite arc midpoints. Their intersection is the $A$-excentre. All preceding algebra still applies, giving $NI_A=|s_1'|^2/2$ and $r_A=||s_1'|^2-1|/2$ for $s_1'=-u+v+w$.

To determine the sign without assuming an excentre-location theorem, put $p=\beta-\alpha>0$, $q=\gamma-\beta>0$, with $p+q<\pi$. Direct trigonometric expansion gives
\[
|s_1'|^2-1
=2(1+\cos p-\cos(p+q)-\cos q)
=8\sin(q/2)\sin((p+q)/2)\cos(p/2)>0.
\]
Therefore $NI_A=1/2+r_A$, proving external tangency. Cyclic relabelling proves the same for the other excircles, including the equilateral triangle.
\end{proof}

\begin{beauty}
Inversion preserves angles and tangencies but changes the quantities that made the original picture awkward. Equal circles appear inside a distorted Steiner ring; a straight stack reveals the Pappus heights. Coefficient vectors then record Descartes' tangencies as fixed pairings and make the replacement rule linear. Their zero-curvature boundary case is a line, a useful reminder that the extended picture contains more than finite circles.
\end{beauty}

\begin{figure}[ht]
\centering
\fig{v9-appb-descartes-line}
\caption{The tangent-circle curvatures $1,1,4,12$ satisfy Descartes' relation. Replacing the curvature-$12$ member gives curvature zero: the common tangent line $y=-1$, not a circle with a finite centre. Exercise~\ref{ex:v9-descartes-line} computes its normalized coefficient vector.}
\label{fig:v9-descartes-line}
\end{figure}

\begin{figure}[ht]
\centering
\fig{v9-appb-image-charge}
\caption{A meridional cross-section of a sphere of radius $R=1$. A source at $P=(2,0)$ has inverse point $P'=(1/2,0)$. At each sphere point $Q$, $|Q-P'|=|Q-P|/2$. A charge of strength $-1/2$ at $P'$ therefore cancels the source potential on the whole sphere; Exercise~\ref{ex:v9-sphere-image-charge} verifies the three-dimensional derivatives and flux.}
\label{fig:v9-sphere-image-charge}
\end{figure}
\FloatBarrier

A circle equation was enough to prove the transformation rules, but the choice of coordinates made the constructions simpler. The same coefficient viewpoint will return in Appendix~C, where projective coordinates let incidence and cross-ratios take the place of Euclidean lengths. Before leaving inversion, the exercises check its finite-circle, line and tangency cases directly.

\section{Exercises}

\begin{exercise}[$\star$]\label{geometry:ex:invline}
Under unit inversion about zero, find the images of $\operatorname{Re}z=1$ and $|z-2|=1$. Give their centres and radii, including the point at infinity where needed to complete the first image.
\end{exercise}
\begin{exercise}[$\star$]\label{geometry:ex:vanschooten}
For equilateral $ABC$, let $P$ lie on the circumcircle arc $BC$ avoiding $A$. Prove $PA=PB+PC$, including the endpoints $P=B,C$.
\end{exercise}
\begin{exercise}[$\star\star$]\label{geometry:ex:twins}
In an arbelos with $AC=2r_1$, $CB=2r_2$, $r_1,r_2>0$, draw the perpendicular to $AB$ at $C$. On each side fit a circle tangent to that perpendicular, the outer semicircle and the small semicircle on that side. Prove that these two circles really lie in the upper arbelos and both have radius $r_1r_2/(r_1+r_2)$.
\end{exercise}
\begin{exercise}[$\star\star$]\label{geometry:ex:steinerformula}
Suppose circles of radii $0<\rho<R$ are strictly nested, with centre distance $d$ satisfying $d+\rho<R$, and $n\ge3$. Prove that a consistently advancing Steiner chain of $n$ circles winding once exists exactly when
\[
(R-\rho)^2-4R\rho\tan^2(\pi/n)=d^2.
\]
\end{exercise}
\begin{exercise}[$\star$]\label{geometry:ex:soddy}
Three externally tangent circles have radii $2,3,10$. Find the signed curvatures and ordinary radii of both circles tangent to all three, distinguishing the enclosing solution.
\end{exercise}
\begin{exercise}[$\star\star$]\label{geometry:ex:pappusradius}
In the arbelos with $AB=1$, $AC=t\in(0,1)$, prove that the advancing Pappus chain has
\[
r_n=\frac{t(1-t)}{2(t+n^2(1-t)^2)},\qquad n\ge0.
\]
Compute its centre height as well and compare with the Pappus theorem.
\end{exercise}
\begin{exercise}[$\star$]\label{geometry:ex:gasketroot}
For root quadruple $(-1,2,2,3)$, list the first two generations of new curvatures, retaining multiplicities or stating which are repeated. Prove by induction that none is divisible by four using the swap rule modulo four.
\end{exercise}
\begin{exercise}[$\star$]\label{geometry:ex:mobius3}
Find the M\"obius map taking $1,i,-1$ to $0,1,\infty$. Determine the image of the open unit disc and justify which side of its boundary is obtained.
\end{exercise}
\begin{exercise}[$\star\star$]\label{geometry:ex:trace}
Let a nonidentity M\"obius map have a determinant-one matrix. Prove that it has one or two fixed points on $\hat\C$. Show that one occurs exactly when $\operatorname{tr}^2=4$, giving conjugacy to $z+1$. With two, prove conjugacy to $\lambda z$, $\lambda\ne0,1$, and
$\lambda+\lambda^{-1}+2=\operatorname{tr}^2$. Deduce the elliptic/loxodromic criterion stated in the remark.
\end{exercise}
\begin{exercise}[$\star\star$]\label{geometry:ex:stereo}
Inversion in a sphere of $\R^3$ sends $x\ne O$ to $O+r^2(x-O)/|x-O|^2$. Prove that spheres and planes map to spheres and planes. Show that inversion about the north pole $N$ of the unit sphere, with radius $\sqrt2$, restricts to stereographic projection onto the equatorial plane. Deduce that circles on the sphere project to circles or lines.
\end{exercise}
\begin{exercise}[$\star$]\label{geometry:ex:orthcircles}
For distinct inverse points $P,Q$ with respect to an ordinary circle $\omega$, prove that every circle through $P,Q$ is orthogonal to $\omega$. Explain the same claim for the generalized circle given by their line.
\end{exercise}
\begin{exercise}[$\star\star\star$]\label{geometry:ex:hexlet}
\emph{Soddy's hexlet and its degenerations.} Three fixed spheres are pairwise tangent at distinct points and bound a common interstitial region, allowing an enclosing sphere. A chain in that region is tangent to all three and advances to the next tangent member in a consistent sense, rather than reversing. Invert at a fixed pair's contact point. Prove six-step closure, allowing a plane as a generalized sphere when inversion requires it. Give the explicit finiteness condition under which all six members are genuine spheres. Why would arbitrary backtracking invalidate the statement?
\end{exercise}

\begin{exercise}[$\star\star$]\label{ex:v9-descartes-line}
\emph{Zero curvature is a line.} Consider tangent circles with centres $-1,1,-3i/4,-5i/12$ and radii $1,1,1/4,1/12$. (a) Verify the six distinct tangencies and compute their normalized vectors. (b) Swap the last member and identify the resulting generalized circle and its three contact points. (c) Compute its coefficient $w=-\bar B$ and explain why replacing that coefficient by zero times a finite centre would give a wrong answer.
\end{exercise}
\begin{exercise}[$\star\star$]\label{ex:v9-disc-mobius}
\emph{Every M\"obius symmetry of the disc.} Prove that $f(D)=D$ for a M\"obius map exactly when
\[
f(z)=e^{i\theta}\frac{z-a}{1-\bar az},\qquad |a|<1.
\]
Verify the disc inequality and inverse formula directly. For the converse, send $a=f^{-1}(0)$ to zero, then use the unit-circle boundary to determine the remaining map. Prove uniqueness of $a$ and $e^{i\theta}$ and compute $|f'(a)|$.
\end{exercise}
\begin{exercise}[$\star\star\star$]\label{ex:v9-sphere-image-charge}
\emph{An image charge for a spherical boundary.} Let $|P|=d>R>0$, $P'=R^2P/d^2$, and
\[
U(x)=\frac1{|x-P|}-\frac{R/d}{|x-P'|}.
\]
(a) Prove $|x-P'|^2-(R^2/d^2)|x-P|^2=(1-R^2/d^2)(|x|^2-R^2)$. Deduce that $U$ vanishes on $|x|=R$ and is positive outside, away from $P$. (b) With $\Delta=\sum_{j=1}^3\partial_j^2$, verify $\Delta U=0$ outside the sphere away from $P$ and $U\to0$ at infinity. (c) Compute its outward radial derivative on the sphere and prove its flux integral is $4\pi R/d$ using the divergence theorem on a punctured ball. This verifies a candidate potential in units where a unit point charge contributes $1/|x-P|$.
\end{exercise}
\begin{exercise}[$\star\star$]\label{ex:v9-ptolemy-defect}
\emph{A quantitative Ptolemy ratio.} For distinct finite complex points $a,b,c,d$, set
\[
z=\frac{(a-b)(c-d)}{(a-c)(b-d)},\qquad
s=\frac{AB\,CD+AD\,BC}{AC\,BD}.
\]
(a) Prove $1-z=(a-d)(b-c)/((a-c)(b-d))$, hence $s=|z|+|1-z|\ge1$. Explain M\"obius invariance whenever the images remain finite and identify equality. (b) For $s>1$, derive
\[
\frac{4(\operatorname{Re}z-1/2)^2}{s^2}
+\frac{4(\operatorname{Im}z)^2}{s^2-1}=1.
\]
(c) If $s\le1+\varepsilon$, $\varepsilon\ge0$, prove $|\operatorname{Im}z|\le\sqrt{2\varepsilon+\varepsilon^2}/2$, including the case $s=1$, and characterize equality in the bound for $\varepsilon>0$.
\end{exercise}

\section{Hints}

\begin{hint}{geometry:ex:invline}
Use the real-axis endpoints and reflection symmetry, or substitute into the generalized circle equation.
\end{hint}
\begin{hint}{geometry:ex:vanschooten}
Apply Ptolemy's equality to the cyclic order $A,B,P,C$, then divide by the common side length.
\end{hint}
\begin{hint}{geometry:ex:twins}
Put $C=0$, $A=-2r_1$, $B=2r_2$. The left twin has centre $(-x,y)$ for radius $x$. Use external tangency to the left small circle and internal tangency to the outer circle, then verify the inequalities selecting the actual upper disc.
\end{hint}
\begin{hint}{geometry:ex:steinerformula}
In the concentric picture $(b-a)/(b+a)=\sin(\pi/n)$. Compute its inversive distance and use the strict nesting to fix the sign.
\end{hint}
\begin{hint}{geometry:ex:soddy}
The curvature roots are $k_1+k_2+k_3\pm2\sqrt{k_1k_2+k_2k_3+k_3k_1}$. An outer circle uses the negative sign of curvature, not a negative radius.
\end{hint}
\begin{hint}{geometry:ex:pappusradius}
Invert about $A=0$ with radius one. The two boundary circles become vertical lines $x=1,1/t$. Invert one circle of the resulting explicit stack back.
\end{hint}
\begin{hint}{geometry:ex:gasketroot}
Do not undo the immediately preceding replacement when listing the next generation. Modulo four, preserve the property of having two entries congruent to two and two odd entries.
\end{hint}
\begin{hint}{geometry:ex:mobius3}
Use the three-point formula and compute the imaginary part for $|z|<1$.
\end{hint}
\begin{hint}{geometry:ex:trace}
For $c\ne0$ the fixed-point quadratic has discriminant $\operatorname{tr}^2-4$. For $c=0$ remember the fixed point at infinity. Normalize one or two fixed points before classifying.
\end{hint}
\begin{hint}{geometry:ex:stereo}
A sphere or plane has equation $A|x|^2+B\cdot x+C=0$. For the projection, substitute $|x-N|^2=2-2x_3$ on the unit sphere.
\end{hint}
\begin{hint}{geometry:ex:orthcircles}
Use the inverse-pair orthogonality proposition; the line $PQ$ passes through the inversion centre.
\end{hint}
\begin{hint}{geometry:ex:hexlet}
The fixed pair becomes parallel planes separated by $h$, and the third sphere becomes a ball of diameter $h$ between them. Chain-ball centres lie in the middle plane, on a circle of radius $h$, with successive distances $h$. A transformed sphere through the inversion centre becomes a plane when brought back.
\end{hint}
\begin{hint}{ex:v9-descartes-line}
Use $X=k(1,-\bar c,|c|^2-r^2)$ and $X_4'=2(X_1+X_2+X_3)-X_4$. Read its line equation before attempting a centre calculation.
\end{hint}
\begin{hint}{ex:v9-disc-mobius}
For $\phi_a=(z-a)/(1-\bar az)$, expand $1-|\phi_a|^2$. A M\"obius map fixing zero has form $\eta z/(1+\beta z)$; constant modulus on $|z|=1$ determines $\beta$.
\end{hint}
\begin{hint}{ex:v9-sphere-image-charge}
For $r=|x|$, $\Delta h(r)=h''(r)+2h'(r)/r$ away from zero, by direct differentiation. On the boundary use the distance ratio from (a). For flux, remove a small ball about $P'$ and track the inward normal of that hole.
\end{hint}
\begin{hint}{ex:v9-ptolemy-defect}
Write $u=|z|$, $v=|1-z|$. Then $u+v=s$ and $(u-v)s=2\operatorname{Re}z-1$. Substitute into $u^2=(\operatorname{Re}z)^2+(\operatorname{Im}z)^2$.
\end{hint}

\section{Solutions}

\begin{solution}{geometry:ex:invline}
\idea Real-axis symmetry and two real endpoints determine each image circle.
\writeup The completed line $\operatorname{Re}z=1$ contains $1$ and infinity, mapping to $1,0$. Its image is a circle symmetric about the real axis, with those diameter endpoints: centre $1/2$, radius $1/2$. Without completing the original line, its finite image omits zero. The circle $|z-2|=1$ has real endpoints $1,3$, mapping to $1,1/3$. Its image circle has centre $2/3$ and radius $1/3$. Direct substitution into the circle equations gives the same results.
\end{solution}

\begin{solution}{geometry:ex:vanschooten}
\idea The equal triangle sides cancel from Ptolemy, leaving the desired sum of distances.
\writeup For $P$ in the open arc, $A,B,P,C$ are in cyclic order. Ptolemy gives $AB\,PC+AC\,BP=AP\,BC$. Divide by the positive common side length to obtain $PC+PB=PA$. At $P=B$ or $C$ it reduces directly to equality of triangle sides, also covered by the theorem's adjacent-coincidence case.
\end{solution}

\begin{solution}{geometry:ex:twins}
\idea Solve the two tangency equations, then check that their solution is the disc required by the geometry rather than an extraneous squared-distance solution.
\writeup Put $C=(0,0)$, with $A=(-2r_1,0)$, $B=(2r_2,0)$. The outer circle has centre $(r_2-r_1,0)$ and radius $R=r_1+r_2$. A left-hand circle of radius $x$ tangent to the vertical line has centre $(-x,y)$. Its two circle tangencies give
\[
(r_1-x)^2+y^2=(r_1+x)^2,\qquad
(r_2-r_1+x)^2+y^2=(r_1+r_2-x)^2.
\]
The first gives $y^2=4r_1x$; substitution in the second yields $4(r_1+r_2)x=4r_1r_2$. Thus $x=r_1r_2/(r_1+r_2)$ and choose $y=2\sqrt{r_1x}>0$.

Here $0<x<r_1$ and $x<R$, so $y>2x>x$: the whole disc lies above the baseline. Its rightmost point has horizontal coordinate zero. The unsquared distances are respectively $r_1+x$ and $R-x>0$, proving external tangency to the left small disc and containment in the outer disc. Its distance to the right small centre has square $(r_2+x)^2+y^2>(r_2+x)^2$, so it does not intrude into that small disc. It really lies in the arbelos and is unique on the upper left. Exchanging $r_1,r_2$ constructs the upper right twin with the same radius.
\end{solution}

\begin{solution}{geometry:ex:steinerformula}
\idea The concentric ratio determines the step angle. Its inversive distance transfers the closing condition to a nonconcentric nested pair.
\writeup In concentric position let the boundary radii be $a<b$. One forward turn with $n\ge3$ circles requires
$(b-a)/(b+a)=\sin(\pi/n)=s\in(0,1)$. Equivalently $b/a=(1+s)/(1-s)$ and
\[
\frac{a^2+b^2}{2ab}
=\frac{1+s^2}{1-s^2}=1+2\tan^2(\pi/n).
\]
The concentric ratio $b/a>1$ is uniquely determined by this value because $(t+t^{-1})/2$ is strictly increasing for $t>1$. For the strictly nested original circles, $\delta=(R^2+\rho^2-d^2)/(2R\rho)>1$, so absolute invariance gives this same positive value. Rearranging yields the stated equation. Conversely that equation gives the required inversive distance, hence the unique concentric ratio and forward step $2\pi/n$. Its equal-circle chain closes after exactly $n$ steps and maps back to the original gap, proving sufficiency.
\end{solution}

\begin{solution}{geometry:ex:soddy}
\idea The two quadratic roots distinguish the small gap from the enclosing circle through their signs.
\writeup The curvature sum is $1/2+1/3+1/10=14/15$, and the pairwise-product sum is $1/6+1/30+1/20=1/4$. Descartes therefore gives $k_4=14/15\pm1$. The positive root $29/15$ has radius $15/29$ in the inner gap. The negative root $-1/15$ is the enclosing circle, whose ordinary radius is $15$. Neither root is zero in this example.
\end{solution}

\begin{solution}{geometry:ex:pappusradius}
\idea The stack has a known centre and radius. A single inverse-circle calculation recovers both the original radius and height.
\writeup Under unit inversion about $A=0$, the boundary circles become $x=1$ and $x=1/t$. Their separation is $w=(1-t)/t>0$. The image of $C_0$ has diameter endpoints $1,1/t$, so the $n$th stack circle has
\[
p=\tfrac12(1+1/t)+inw,\qquad s=w/2.
\]
For an ordinary circle $|z-p|=s$ avoiding zero, substitution $z=1/\bar\zeta$ and completion of the square give image centre $p/(|p|^2-s^2)$ and radius $s/\bigl||p|^2-s^2\bigr|$. Here the denominator is positive:
\[
D=|p|^2-s^2=\frac1t+\frac{n^2(1-t)^2}{t^2}.
\]
Thus
\[
r_n=\frac{s}{D}
=\frac{t(1-t)}{2(t+n^2(1-t)^2)},\qquad
\text{height}=\frac{nw}{D}=2nr_n.
\]
The last quantity is $n$ diameters. At $n=0$ the formula gives the boundary circle radius $(1-t)/2$ and zero height, as required.
\end{solution}

\begin{solution}{geometry:ex:gasketroot}
\idea The swap is an integer linear operation. A preserved residue pattern proves the exclusion at every generation, rather than just in a finite list.
\writeup Replacing each member of $(-1,2,2,3)$ gives first-generation curvatures $15,6,6,3$. For the second generation, swap any of the other three members of each new quadruple, never immediately undoing the latest swap. The lists from the four branches are respectively
\[
(38,38,35),\quad(23,14,11),\quad(23,14,11),\quad(15,6,6).
\]
Their distinct values are $6,11,14,15,23,35,38$; the displayed branches retain all twelve multiplicities.

Modulo four, the root has two entries congruent to two and two odd entries. Replacing an odd entry $k$ leaves an odd sum among the other three, so its new residue is $2-k$, still odd. Replacing an entry congruent to two leaves an even sum among the others, so its new residue is $-2=2$ modulo four. The pattern persists under every swap by induction. No entry can be zero modulo four, proving the claim for every generation.
\end{solution}

\begin{solution}{geometry:ex:mobius3}
\idea Normalize the three boundary points, then use an explicit imaginary-part formula to select the image component.
\writeup The triple formula gives
\[
f(z)=\frac{(z-1)(i+1)}{(z+1)(i-1)}
=-i\frac{z-1}{z+1}.
\]
It sends $1,i,-1$ to $0,1,\infty$, respectively, and
$\operatorname{Im}f(z)=(1-|z|^2)/|z+1|^2$. Hence it sends the disc into the upper half-plane. Its inverse is $z=(i-w)/(i+w)$; for $\operatorname{Im}w>0$, $|i-w|<|i+w|$, so the inverse lies in the disc. The image is the entire upper half-plane.
\end{solution}

\begin{solution}{geometry:ex:trace}
\idea Fixed points can be moved to infinity or to zero and infinity. The normal forms then expose the trace invariant.
\writeup For $c\ne0$, finite fixed points solve $cz^2+(d-a)z-b=0$, whose discriminant is $(a+d)^2-4$ because $ad-bc=1$. Infinity is not fixed in this case. For $c=0$, infinity is fixed and the map is $(a/d)z+b/d$, with $ad=1$. If $a\ne d$, it has one additional finite fixed point. If $a=d$, then $a=d=\pm1$, and nonidentity forces $b\ne0$, giving only infinity. Thus there is one fixed point exactly for squared trace four.

Conjugation by a M\"obius matrix preserves trace before determinant normalization; choosing the other determinant-one sign changes only its sign. Squared trace is therefore invariant. Move a unique fixed point to infinity. The affine normal form must then have multiplier one, otherwise it has a finite fixed point too, so it is $z+\beta$, $\beta\ne0$. Scaling by $1/\beta$ makes it $z+1$. For two fixed points, move them to zero and infinity; the map becomes $\lambda z$, with $\lambda\ne0,1$. A determinant-one representative is $\operatorname{diag}(\lambda^{1/2},\lambda^{-1/2})$, giving $\operatorname{tr}^2=\lambda+\lambda^{-1}+2$.

If $|\lambda|=1$ and $\lambda\ne1$, this expression is $4\cos^2(\theta/2)\in[0,4)$ for $\lambda=e^{i\theta}$, the elliptic case. Conversely a value in that interval makes the two reciprocal roots of $\lambda+\lambda^{-1}=\operatorname{tr}^2-2$ unit complex numbers different from one. All remaining two-fixed-point cases have $|\lambda|\ne1$ and are loxodromic. The identity, excluded throughout, has infinitely many fixed points and must not be put into the nonidentity trichotomy.
\end{solution}

\begin{solution}{geometry:ex:stereo}
\idea The same quadratic equation describes spheres and planes in three dimensions. Inversion exchanges its constant and quadratic coefficients.
\writeup For unit inversion at zero, $y=x/|x|^2$ and $x=y/|y|^2$. The equation $A|x|^2+B\cdot x+C=0$ becomes
$C|y|^2+B\cdot y+A=0$ after multiplying by $|y|^2$. Nondegenerate spheres and planes therefore map to spheres or planes, with the usual zero/infinity completions.

For $N=(0,0,1)$ and inversion radius $\sqrt2$,
\[
\mathcal I(x)=N+\frac{2(x-N)}{|x-N|^2}.
\]
On $|x|=1$, $x\ne N$, the denominator is $2-2x_3$, so the third coordinate is zero and the first two are $x'/(1-x_3)$. This is exactly the stereographic formula from Chapter~1.

A sphere circle is the intersection of the unit sphere and a plane. Its image is the equatorial plane intersected with the image of that plane, hence a circle or line. It cannot collapse to a point or be empty, by injectivity; it cannot be the entire equatorial plane, since the original plane was not the unit sphere. If the original circle contains $N$, its image contains infinity and is a line; otherwise it is an ordinary circle.
\end{solution}

\begin{solution}{geometry:ex:orthcircles}
\idea A distinct inverse pair fixes the power of the inversion centre relative to every circle through the pair.
\writeup For any circle through $P,Q=P'$, that power is $OP\,OQ=r^2$, so its centre distance satisfies $OC^2=\rho^2+r^2$. The right-triangle criterion proves orthogonality, exactly as in Proposition~\ref{geometry:prop:orthfixed}. The line $PQ$ passes through $O$; it meets the inversion circle along a diameter and is perpendicular to its tangent at both endpoints, so it is orthogonal in the generalized-circle sense too.
\end{solution}

\begin{solution}{geometry:ex:hexlet}
\idea The three-dimensional inversion replaces a curved gap by the space between two planes outside one ball. Centres of all moving balls then lie on a circle with a fixed $60^\circ$ step.
\writeup Invert at the contact point of two fixed spheres. Their common tangent plane gives two parallel image planes, separated by some $h>0$. Distinct fixed tangencies make the third sphere avoid the inversion centre, so its image is an ordinary sphere tangent to both planes. Its centre lies in their middle plane and its radius is $h/2$: equality of the two distances to distinct parallel planes forces that centre and radius. The proper common gap maps to the region between the planes and outside that third ball, also when the original configuration has an enclosing member.

Each moving generalized sphere becomes a ball tangent to both planes and externally tangent to the third ball. It has radius $h/2$ and centre in the middle plane at distance $h$ from the fixed centre. Two consecutive moving balls are tangent exactly when their centres have distance $h$. On the radius-$h$ circle, this chord subtends $60^\circ$. Choosing one consistent advancing sense therefore gives six distinct centres before the first return, for every starting centre. Invert back to obtain six-step closure.

A transformed moving ball becomes a plane precisely when its boundary contains the inversion centre used in bringing the configuration back. Thus all six original members are genuine finite spheres exactly when none of those six transformed boundaries passes through that point. Otherwise closure is still valid in generalized spheres, with the corresponding plane members; an assertion about six finite spheres would need that extra finiteness condition. Immediate reversal onto the previous centre instead gives a two-step retracing sequence, not an advancing hexlet. This is why the progression and degeneration conventions belong to the statement.
\end{solution}

\begin{solution}{ex:v9-descartes-line}
\idea Read the replacement from its coefficient equation. A zero first coefficient removes the quadratic term and leaves a line, not an ordinary centre-radius formula.
\writeup (a) The first two centre distances are two. Their distances to $-3i/4$ are $5/4=1+1/4$, and to $-5i/12$ are $13/12=1+1/12$. The last two centres are distance $1/3=1/4+1/12$ apart. Thus every pair is externally tangent. For three of these discs to share a contact point, their centres would have to lie on one line through that point; the listed centre positions show that no triple does. Hence the six contacts are distinct. Their vectors are
\[
X_1=(1,1,0),\quad X_2=(1,-1,0),\quad
X_3=(4,-3i,2),\quad X_4=(12,-5i,2).
\]

(b) The swap gives $X_4'=(0,-i,2)$. Its equation is $-iz+i\bar z+2=0$, or $y=-1$. The first two unit circles touch it at $(-1,-1),(1,-1)$, and the radius-$1/4$ circle touches it at $(0,-1)$. These are three different contacts. The curvature equation gives $k_4'=2(1+1+4)-12=0$, agreeing with the line.

(c) Its weighted coefficient is $w_4'=-\overline{-i}=-i$, and the vector swap gives the same value from $2(-1+1-3i)-(-5i)$. Replacing it by zero times a finite centre would incorrectly make it zero. A line has no finite circle centre; its nonzero linear coefficient retains the direction information that such a product would erase.
\end{solution}

\begin{solution}{ex:v9-disc-mobius}
\idea Move the unique preimage of zero to zero. A boundary-preserving M\"obius map fixing zero then has no remaining denominator term.
\writeup Let $\phi_a(z)=(z-a)/(1-\bar az)$. Expansion gives
\[
1-|\phi_a(z)|^2=
\frac{(1-|a|^2)(1-|z|^2)}{|1-\bar az|^2}.
\]
For $|a|<1$ the denominator has no zero on the closed disc, and this sends the open disc into itself and its boundary onto the boundary. The inverse is $(w+a)/(1+\bar aw)$, which satisfies the analogous inequality, so $\phi_a(D)=D$. Multiplication by a unit complex number preserves the disc too.

Conversely let $f(D)=D$. Its inverse sends zero to a unique $a\in D$. The map $h=f\phi_a^{-1}$ preserves $D$ and fixes zero. As a homeomorphism of the extended plane it maps the boundary of $D$ onto itself. Write $h(z)=\eta z/(1+\beta z)$ with $\eta\ne0$, using $h(0)=0$. There is no boundary pole. On $|z|=1$ we have $|1+\beta z|^2=|\eta|^2$. Expanding gives $1+|\beta|^2+2\operatorname{Re}(\beta z)$ constant for all unit $z$, which forces $\beta=0$ (compare $z=1,-1,i,-i$). Hence $|\eta|=1$, giving the stated form.

The unique zero of $f$ determines $a$, and $f\phi_a^{-1}$ determines $e^{i\theta}$, proving uniqueness of both parameters. Differentiation gives
$\phi_a'(z)=(1-|a|^2)/(1-\bar az)^2$, so $|f'(a)|=1/(1-|a|^2)$.
\end{solution}

\begin{solution}{ex:v9-sphere-image-charge}
\idea Inversion gives the boundary distance ratio, so the charge strengths cancel there. Ordinary differentiation and the divergence theorem verify the field properties without a general boundary-value theorem.
\writeup (a) Substitute $P'=R^2P/d^2$ and expand both squared distances. Their cross terms cancel, leaving
\[
|x-P'|^2-\frac{R^2}{d^2}|x-P|^2
=\left(1-\frac{R^2}{d^2}\right)(|x|^2-R^2).
\]
On the boundary the distances are in ratio $R/d$, making $U=0$. Outside it the right side is positive, so $|x-P'|>(R/d)|x-P|$ and $U>0$ away from $P$.

(b) For a radial function $h(r)$, differentiating $r=|x|$ gives $\Delta h=h''+2h'/r$ for $r>0$. For $h(r)=1/r$ this is $2/r^3-2/r^3=0$. Translation gives harmonicity of each source term off its pole. The image point lies inside $|x|<R$, so $\Delta U=0$ in the exterior away from $P$. Both inverse distances tend to zero at infinity, giving the stated limit.

(c) Write $x=Rn$, $|n|=1$, and $t=P\cdot n$. Differentiating in the radial direction and using the boundary ratio gives
\[
\begin{aligned}
\partial_rU(Rn)
&=-\frac{R-t}{|Rn-P|^3}
+\frac{R}{d}\,
\frac{R-R^2t/d^2}{(R/d)^3|Rn-P|^3}\\
&=\frac{d^2-R^2}{R|Rn-P|^3}.
\end{aligned}
\]
For its integral, apply the divergence theorem from Chapter~3 to $\nabla U$ in the ball of radius $R$ with a small radius-$\varepsilon$ ball about $P'$ removed. It is smooth there and has zero divergence. The domain's normal on the small hole points \emph{towards} $P'$. The image term has gradient $(R/d)(x-P')/\varepsilon^3$ on that small sphere, so its flux with this inward normal is $-4\pi R/d$. The other term is smooth near $P'$ and its small-sphere flux tends to zero. Therefore zero total flux gives
\[
\int_{|x|=R}\partial_rU\,\dd S=\frac{4\pi R}{d}.
\]
The small ball lies strictly inside the large one for sufficiently small $\varepsilon$, so every boundary and orientation used here is legitimate.
\end{solution}

\begin{solution}{ex:v9-ptolemy-defect}
\idea Normalize the four-point distance expression into the sum of two complex lengths. The triangle inequality gives equality, and its two-distance algebra gives the quantitative bound.
\writeup (a) Expanding the numerator difference gives
\[
(a-c)(b-d)-(a-b)(c-d)=(a-d)(b-c).
\]
Taking absolute values proves $s=|z|+|1-z|\ge1$. The number $z$ is the reordered cross-ratio $[a,d;b,c]$, hence invariant under M\"obius maps; the expression for $s$ makes it invariant whenever its distance formula remains finite. Equality in $z+(1-z)=1$ requires $z$ real between zero and one. With four distinct points neither endpoint can occur, so $0<z<1$. This is the ordered generalized-circle case of Ptolemy.

(b) Put $x=\operatorname{Re}z$, $y=\operatorname{Im}z$, $u=|z|$, $v=|1-z|$. Then $u+v=s$ and $u^2-v^2=2x-1$, giving $u=(s+(2x-1)/s)/2$. Substitute into $u^2=x^2+y^2$ and simplify to
\[
4y^2=(s^2-1)\left(1-\frac{(2x-1)^2}{s^2}\right).
\]
For $s>1$ division gives the required identity, obtained without assuming any classification of conics.

(c) It follows that $|y|\le\sqrt{s^2-1}/2$, with equality for $s>1$ exactly when $x=1/2$. Since $s\le1+\varepsilon$, the stated bound follows. For $s=1$, part (a) gives $y=0$ directly, so there is no division by zero. For $\varepsilon>0$, equality in the final bound requires both $s=1+\varepsilon$ and $x=1/2$, and then $y=\pm\sqrt{2\varepsilon+\varepsilon^2}/2$. This measures the normalized complex ratio's distance from the real line, rather than an unscaled Euclidean displacement of the original points.
\end{solution}

\section*{Sources and further reading}

\begin{itemize}
\item R.~Descartes, letter to Princess Elisabeth of November 1643, in \emph{\OE uvres de Descartes}, ed.\ Adam and Tannery, vol.~4; F.~Soddy,
\emph{The Kiss Precise}, Nature 137 (1936), 1021 (the epigraph is its second stanza).
\item J.~C.~Lagarias, C.~L.~Mallows and A.~R.~Wilks, \emph{Beyond the Descartes circle theorem}, American Mathematical Monthly 109 (2002), 338--361
(for the coefficient-vector proof and its centre identities).
\item R.~L.~Graham, J.~C.~Lagarias, C.~L.~Mallows, A.~R.~Wilks and C.~H.~Yan, \emph{Apollonian circle packings: number theory}, Journal of Number
Theory 100 (2003), 1--45; J.~Bourgain and A.~Kontorovich, \emph{On the local-global conjecture for integral Apollonian gaskets}, Inventiones
Mathematicae 196 (2014), 589--650; S.~Haag, C.~Kertzer, J.~Rickards and K.~E.~Stange, \emph{The local-global conjecture for Apollonian circle
packings is false}, Annals of Mathematics 200 (2024).
\item C.~T.~McMullen, \emph{Hausdorff dimension and conformal dynamics III: computation of dimension}, American Journal of Mathematics 120 (1998),
691--721 (for the quoted gasket-dimension computation).
\item D.~Pedoe, \emph{Geometry: A Comprehensive Course}, Dover, 1988 (inversion, inversive distance, Steiner's porism); H.~S.~M.~Coxeter, \emph{The
inversive plane and hyperbolic space}, Abhandlungen aus dem Mathematischen Seminar der Universit\"at Hamburg 29 (1966), 217--242.
\item H.~S.~M.~Coxeter and S.~L.~Greitzer, \emph{Geometry Revisited}, MAA, 1967, Chapter~5 (inversion, Feuerbach by inversion, Steiner's porism).
\item T.~Needham, \emph{Visual Complex Analysis}, Oxford University Press, 1997, Chapter~3 (for geometric drawings of the M\"obius and inversion arguments).
\item L.~Bankoff, \emph{Are the twin circles of Archimedes really twins?}, Mathematics Magazine 47 (1974), 214--218; Pappus, \emph{Collection},
Book~IV (the chain in the arbelos).
\item W.~Thomson, \emph{Extrait d'une lettre de M.~William Thomson \`a M.~Liouville}, Journal de Math\'ematiques Pures et Appliqu\'ees 10 (1845),
364--367 (inversion in electrostatics, the method of images).
\end{itemize}
